\begin{abstract}
We propose a novel memory module for building video generators capable of interactively exploring environments.
Previous approaches have achieved similar results either by out-painting 2D views of a scene while incrementally reconstructing its 3D geometry---which quickly accumulates errors---or by using video generators with a short context window, which struggle to maintain scene coherence over the long term.
To address these limitations, we introduce \textit{\mnamel} (\textit{\mname}), a memory module that remembers past views by indexing them geometrically based on the 3D surface elements (surfels) they have observed.
\mname{} enables efficient retrieval of the most relevant past views when generating new ones.
By focusing only on these relevant views, our method produces consistent explorations of imagined environments at a fraction of the computational cost required to use all past views as context.
We evaluate our approach on challenging long-term scene synthesis benchmarks and demonstrate superior performance compared to existing methods in maintaining scene coherence and camera control.
\end{abstract}
